\documentclass[a4paper,11pt]{article}
\usepackage[utf8]{inputenc}
\usepackage[english]{babel}
\usepackage[margin=2.5cm]{geometry}
\usepackage{amsmath}
\usepackage{amssymb}
\usepackage{graphicx}
\usepackage{xcolor}
\usepackage{enumitem}
\usepackage{booktabs}

\title{\textbf{TAF ISC\\UE « Architecture et Ingénierie des Systèmes de Transmission »\\Examen partie optique}}
\author{Décembre 2022}
\date{}

\begin{document}

\maketitle

\section*{Problem context and data}

Two coastal countries denoted A and B are linked by three already deployed optical submarine cables as shown in figure 1.

\begin{figure}[h]
    \centering
    \textit{Figure 1: the three submarine routes linking A and B}
\end{figure}

The main characteristics of the three existing optical routes are summarized in Table 1 below. You cannot modify any of these.

\begin{table}[h]
\centering
\begin{tabular}{lccc}
\toprule
& \textbf{Northern route} & \textbf{Central route} & \textbf{Southern route} \\
\midrule
Length (km) & 3500 & 3000 & 4000 \\
Lineic attenuation $\alpha$ (dB/km) & 0.17 & 0.18 & 0.16 \\
Amplifier spacing (km) & 65 & 50 & 75 \\
EDFA bandwidth (GHz) & 4000 & 3800 & 4200 \\
EDFA NF (dB) & 5.5 & 6 & 5 \\
Max EDFA output power (dBm) & 15 & 15 & 16 \\
\bottomrule
\end{tabular}
\caption{Main characteristics of the three existing submarine cable routes}
\end{table}

As current systems are becoming obsolete, you have to upgrade the three links to coherent WDM transmission while maximizing the link capacity.

The proposed configurations must have at least 1 dB OSNR margin.

The transponders have to be chosen among the ones presented in Table 2:

\begin{table}[h]
\centering
\begin{tabular}{lcccc}
\toprule
\textbf{Model} & \textbf{1} & \textbf{2} & \textbf{3} & \textbf{5} \\
\midrule
Symbol rate & 32 Gbaud & 32 Gbaud & 45 Gbaud & 64 Gbaud \\
Modulation format & DP-QPSK & DP-16QAM & DP-8QAM & DP-16QAM \\
Payload bit rate & 100 Gbit/s & 200 Gbit/s & 200 Gbit/s & 400 Gbit/s \\
Spectral slot occupancy (note 1) & 3 & 3 & 5 & 6 \\
Minimum OSNR (note 2) & 12 dB & 19 dB & 15 dB & 25 dB \\
\bottomrule
\end{tabular}
\caption{Main characteristics of the available coherent transponders}
\end{table}

\textbf{Note 1:} The WDM spectrum is divided in 12.5 GHz wide slots. An optical channel in the WDM multiplex occupies a set of contiguous slots.

\textbf{Note 2:} OSNR is measured and calculated with an optical bandwidth of 12.5 GHz.

\section*{Questions and replies}

\textbf{Name:} \rule{8cm}{0.4pt}

\begin{enumerate}
    \item Write down the Shannon-Hartley formula that gives the theoretical capacity (in bit/s) of a communication channel with additive Gaussian noise:
    
    \vspace{2cm}
    
    \item Write down the formula you use to calculate the OSNR value at the end of a regularly optically amplified fiber line:
    
    \vspace{2cm}
    
    \item Write down the formula you use to calculate the line rate of an optical transponder with polarization multiplexing:
    
    \vspace{2cm}
    
    \item Which parameter do you calculate to select the most suitable transponder(s)?
    
    \vspace{2cm}
    
    \item Fill in the table below (per fiber values):
    
    \begin{table}[h]
    \centering
    \begin{tabular}{lccc}
    \toprule
    & \textbf{Northern route} & \textbf{Central route} & \textbf{Southern route} \\
    \midrule
    Spectral slot count & & & \\
    Transponder model & & & \\
    Channel count & & & \\
    Power per channel (dBm) & & & \\
    End of line OSNR (dB) & & & \\
    Theoretical capacity (Tbit/s) & & & \\
    Actual capacity (Tbit/s) & & & \\
    \bottomrule
    \end{tabular}
    \end{table}
    
    \item What could possibly be done to improve the actual capacity of the links?
    
    \vspace{3cm}
\end{enumerate}

\end{document}
